\documentclass[12pt]{article}
\usepackage[T1]{fontenc}
\usepackage[utf8]{inputenc}
\usepackage{lmodern}
\usepackage{textcomp}
\usepackage{enumitem}
\usepackage{geometry}
\geometry{margin=1in}
\begin{document}
\begin{center}
Answer Key - History - K-12 Level Assessment 2 \
\small (Answers are ordered exactly as in the question paper.)
\end{center}

\section*{Multiple Choice}
\begin{enumerate}[label=\arabic*.]
\item \textbf{Answer:} C
\\ \textit{Options:} A) Aurignacian; B) Acheulean; C) Mousterian; D) both b and c
\\ \textit{Note:} The Mousterian tool technology is specifically associated with Neandertals, characterized by its use of flint and diverse tool types for hunting and processing food, reflecting their adaptation to their environment during the Middle Paleolithic period.
\item \textbf{Answer:} A
\\ \textit{Options:} A) Alaska.; B) Mexico.; C) Canada.; D) the continental United States.
\\ \textit{Note:} Clovis points, which are distinctive stone tools associated with prehistoric North American hunters, have primarily been found in the continental United States and parts of Mexico, but there is no archaeological evidence of their presence in Alaska.
\item \textbf{Answer:} D
\\ \textit{Options:} A) agriculture.; B) irrigation technology.; C) social stratification.; D) a food surplus.
\\ \textit{Note:} A food surplus is essential for the development of social complexity because it allows for population growth, specialization of labor, and the establishment of social hierarchies, which are crucial for more advanced societal structures.
\item \textbf{Answer:} C
\\ \textit{Options:} A) stone tablets.; B) pyramids.; C) monumental works.; D) irrigation canals.
\\ \textit{Note:} Monumental works, such as pyramids, temples, and fortifications, serve as physical representations of the power, resources, and organizational capabilities of ancient state societies, making them the most obvious material symbols of these civilizations.
\item \textbf{Answer:} A
\\ \textit{Options:} A) fruits, leaves, and C$_{3}$ plants.; B) meat derived from carnivorous mammals.; C) meat, fish, and some plants.; D) C$_{4}$ grasses that may have included wild wheat.
\\ \textit{Note:} The correct answer is supported by archaeological evidence showing that early hominids, such as those discovered in Malapa Cave, primarily consumed a diet rich in fruits, leaves, and C$_{3}$ plants, which are indicative of their adaptation to a forested environment.
\end{enumerate}

\section*{True / False}
\begin{enumerate}[label=\arabic*.]
\setcounter{enumi}{5}
\item \textbf{Answer:} False
\\ \textit{Options:} A) True; B) False
\\ \textit{Note:} The statement is false because the first pharaohs of Egypt emerged around 3100 BCE, which is approximately 5100 years ago, not 4100 years ago.
\item \textbf{Answer:} False
\\ \textit{Options:} A) True; B) False
\\ \textit{Note:} The statement is false because while food surpluses, a standing army, artwork, and a system of plumbing can be characteristics of advanced civilizations, they are not universally required for a civilization to exist, as some early societies may not have developed all these features.
\item \textbf{Answer:} False
\\ \textit{Options:} A) True; B) False
\\ \textit{Note:} While hieroglyphics is one of the earliest known writing systems, cuneiform, developed by the Sumerians in Mesopotamia around 3200 BCE, predates hieroglyphics and is considered the oldest form of writing.
\item \textbf{Answer:} True
\\ \textit{Options:} A) True; B) False
\\ \textit{Note:} By$_{1900}$B.P., the Mayans primarily used stone tools and did not have access to bronze metallurgy, which was not developed in their civilization.
\item \textbf{Answer:} True
\\ \textit{Options:} A) True; B) False
\\ \textit{Note:} Maize, beans, and squash, often referred to as the "Three Sisters," formed the foundation of agriculture for many indigenous New World civilizations, providing essential nutrients and supporting sustainable farming practices that enabled these societies to thrive.
\end{enumerate}

\section*{Long Form Response}
\begin{enumerate}[label=\arabic*.]
\setcounter{enumi}{10}
\item \textbf{Answer:} The Mississippian culture, which thrived in the southeastern United States from around 800 CE to 1600 CE, was characterized by advanced agricultural practices that heavily relied on the cultivation of maize and squash. These crops were not only staples in their diet but also supported large populations and complex societies. The ability to grow maize allowed the Mississippians to settle in one place, leading to the development of large, organized communities and the construction of impressive earthen mounds. Additionally, the cultivation of squash provided essential nutrients and complemented their agricultural system, facilitating trade and social interactions among different groups. Thus, the reliance on these crops significantly shaped their society, influencing their social structure, economy, and cultural practices.
\\ \textit{Note:} correct because it effectively outlines how the cultivation of maize and squash enabled the Mississippian culture to establish permanent settlements, which in turn fostered the growth of large, organized societies and complex social structures.
\item \textbf{Answer:} The evidence supporting trade between the Indus Valley cities and the city-states of ancient Mesopotamia is notably highlighted by the discovery of Harappan seals in Mesopotamia and Mesopotamian cylinder seals in the Indus Valley. These seals, which often feature intricate designs and symbols, suggest that both civilizations engaged in commerce, exchanging goods and possibly ideas. The presence of these seals indicates not only a direct economic relationship but also cultural interactions, as they were used for marking ownership and authenticity of traded items. This exchange likely facilitated the movement of goods such as textiles, metals, and agricultural products, enhancing the prosperity of both regions. Overall, the archaeological findings of seals provide compelling evidence of a sophisticated trading network linking the Indus Valley and Mesopotamia.
\\ \textit{Note:} correct because it identifies the discovery of Harappan and Mesopotamian seals as key evidence of trade and cultural exchange between the two ancient civilizations, illustrating both economic interactions and the sharing of ideas.
\end{enumerate}

\vfill
\noindent\textit{End of Answer Key}
\end{document}